\section{Training Details}
\label{sec:appendix_training}

\subsection{Phase~1: Behavior Cloning and Phase~2: REINFORCE}

The Neural GibbsQ pipeline has two sequential phases. The initial BC phase trains the neural policy by supervised imitation of an analytical teacher~\cite{ross2011reduction}, minimizing cross-entropy between the network output and the teacher's routing distribution. Starting from the BC-pretrained weights, we then apply \reinforce{}~\cite{williams1992simple} to optimize steady-state mean total queue length. Table~\ref{tab:unified_training_config} summarizes the configuration and ablation variants.

\begin{table}[htbp]
\centering
\small
\caption{Consolidated configuration for Neural Training (BC + REINFORCE) and Ablation Variants.}
\label{tab:unified_training_config}
\renewcommand{\arraystretch}{1.25}
\begin{tabular}{@{\hspace{6pt}}l @{\hspace{20pt}} p{8.5cm}@{\hspace{6pt}}}
\toprule
\textbf{Phase 1: Behavior Cloning} & \textbf{Value / Description} \\
\midrule
Expert policy         & UAS ($\temp = 10.0$); Reflected UAS for reflected-teacher \\
Preprocessing         & $\log(1 + (\Qlen{i}+1)/\srvrate_i)$ (look-ahead potential) \\
Training steps        & 1{,}000 \\
Final accuracy        & 99.2\% (best seed, 97--99\% inter-seed range) \\
Loss / Optimizer      & Cross-entropy / AdamW ($\alpha = 2 \times 10^{-3}$, weight decay $10^{-4}$) \\
\midrule
\textbf{Phase 2: REINFORCE} & \textbf{Value / Description} \\
\midrule
Initialization        & BC pretrained weights \\
Ensemble seeds        & 5 independent initializations (select best on validation) \\
Batch size / Epochs   & 8 episodes / 8 (ablation budget) \\
Sim.\ time / episode  & 750 time units (ablation budget) \\
Optimizer             & Adam ($\alpha = 1 \times 10^{-4}$) \\
Baseline              & Running average of episode returns \\
\midrule
\textbf{Ablation Variants} & \textbf{Description and Rationale} \\
\midrule
BC: UAS $\to$ \reinforce{} & Default pipeline. Tests standard two-phase training. \\
\addlinespace[3pt]
BC: Cal.\ UAS $\to$ \reinforce{} & Replaces BC teacher. Isolates the teacher-quality effect. \\
\addlinespace[3pt]
No Log-Norm           & Removes $\log(1 + (\Qlen{i}+1)/\srvrate_i)$ preprocessing. Tests input stability. \\
\addlinespace[3pt]
Zero-Init Final       & Zero-initializes final layer weights before BC training. \\
\addlinespace[3pt]
\reinforce{} from Scratch & Skips BC entirely. Tests if \reinforce{} alone converges. \\
\bottomrule
\end{tabular}
\end{table}

The exact same network architecture, preprocessing, learning rate, and optimizer settings are applied to the reflected-teacher ablation variant, yielding comparable target fidelity. For \reinforce{}, gradient clipping (global norm 0.5) and a small entropy regularization bonus ($0.01$) are applied during training for stability; the BC initialization provides the primary safeguard against unstable updates. Convergence is typically observed by epoch 6--8 under the ablation budget, verifying 8 epochs as sufficient for this architecture. The 32.6\% magnitude error relative to finite-difference gradients (Appendix~\ref{sec:appendix_gradient}) effectively reduces the learning rate but does not affect the optimization direction.

\begin{figure}[htbp!]
\centering
\includegraphics[width=0.75\textwidth, height=0.35\textheight, keepaspectratio]{reinforce_training_curve.png}
\caption{\reinforce{} training curve: deterministic policy score vs.\ epoch.
  Rapid convergence from the BC initialization is visible; the score plateaus by epoch 6--8.}
\label{fig:reinforce_training}
\end{figure}

The main ablation result is that teacher choice ($1.32$ units) has a much larger effect than the tested architectural modifications (all below $0.2$ units), while training from scratch without BC performs very poorly (mean queue $> 5{,}000$).

\paragraph{Training--evaluation horizon mismatch}
The ablation training episodes use a horizon of 750 time units, while the evaluation protocol (Protocol~B) uses 20{,}000 time units. This 26.7$\times$ disparity may introduce a distributional mismatch between the training and evaluation regimes. However, we note that effective transient recovery---rapidly clearing queue buildups---is precisely the behavior that determines steady-state quality in stable systems. The empirical consistency between short-horizon training scores and long-horizon evaluation metrics (Section~\ref{sec:experiments}) suggests the mismatch does not systematically degrade performance in this setting.

\paragraph{Compute budget}
Wall-clock training times were not systematically logged; the total compute budget for the full experimental campaign is approximately 200 CPU-hours on Intel Xeon E5-2680 v4 (2.4\,GHz, 28 cores). Approximate wall-clock times per experiment: drift verification $\sim$2\,min, policy benchmark $\sim$4\,hr, ablation $\sim$12\,hr, generalization $\sim$8\,hr.

\section{Piecewise Lyapunov Candidate Search}
\label{sec:appendix_piecewise_search}

The piecewise Lyapunov candidate presented in Section~\ref{sec:caluas_piecewise_candidate} and \cref{eq:piecewise_candidate} was numerically discovered using an automated post-validation search procedure. The search space consisted of convex combinations of four template components:
\begin{enumerate}
\item a heterogeneous quadratic term $\frac{1}{2}\sum_i \Qlen{i}^2 / \mu_i$,
\item an unweighted global-load term $(\sum_i \Qlen{i})^2$,
\item a cubic heterogeneous term $\sum_i \Qlen{i}^3 / \mu_i$, and
\item a maximum-queue term $\max_i \Qlen{i}^2$.
\end{enumerate}

The scalar coefficients were optimized with L-BFGS-B to reduce the largest observed value of the exact generator $\gen \lyap_{\mathrm{PW}}(\Qstate)$ over the shell family $|\Qstate|_1 \in \{12, \dots, 20\}$. States that repeatedly appeared as positive-drift counterexamples in auxiliary random-shell checks were up-weighted in later search rounds. The final coefficient vector was then fixed at six-decimal precision for the verification capsule reported in Appendix~\ref{sec:appendix_piecewise}.
